\section{Conclusion}
% ================================================================================

This document has specified the Bridge Contract, a tuple
$BC = (S, E, R, C)$ that fixes the obligations of a bridge component
during client detach, human unavailability, and session
resume. The contract specifies one class of the obligations that
AgentThread \cite{zheng2026formal} reports unassigned across protocols,
here at the intersection of AGTP's transport scope, GAR's audit scope,
and PACE's resumption scope; it does not show that class to be complete. The contract is
positioned as a complement to---not a replacement for---AGTP, GAR,
and PACE: AGTP defines the transport headers the bridge populates;
GAR defines the audit record shape the bridge emits into (under the
GEC-less deployment profile of Definition~\ref{def:gecless} when no
GEC is present); PACE defines the re-attach mechanism for which the
bridge keeps a resume record. The contract comprises a
product state machine (Attachment $\times$ Human $\times$ Lifecycle, with
the state names kept only as a derived view), core and extended
conformance invariants I1--I15 including a challenge-response re-attach
with proof of possession over a single-use nonce, answer binding, several
clients per principal, durable append and crash recovery, a statement of
what the bridge can see, and a safety argument, explicit axioms and lemmas,
liveness properties dependent on derived timeouts, and a conformance test
suite of executable tests (example-based, property-based, and
deterministic edge-case).

The criterion of validity is independent implementability: if two
developers, reading this specification, can produce bridges whose logs
agree on the same input trace --- the same committed event kinds and the
same final cell for the vectors V01--V134 --- the specification is
independently implementable (relative to the vectors V01--V134). The Bridge Contract does not claim on-the-wire
interoperability between bridges (the suite is a self-test, not an
interop harness; see Section~\ref{sec:limitations}). The evidence
delivered with this document falls short of that criterion and supports
a weaker claim. TLC finds no invariant violation in the product model,
key lifecycle included, on bounded configurations (two sessions $\times$ two principals
$\times$ two permissions, and a deeper single-session run). I4, I5 and I14
are flag abstractions in the model, the I6 state invariant and the I11(f)
no-eviction property are not meaningfully checked as configured, and I8 is
checked by tests only (see the limitations register). The
two-session run never fires Disconnect, Resolve, Submit or Drain, so
the detach paths are checked only in single-session runs. The model
also rejects every one of its deliberately broken variants. The Python
reference bridge passes its tests, and each of the invariants I1--I15
has a mutant that the suite kills, and no mutant survives. The
language-neutral vectors pass on the Python bridge and (except V52--V54, V88 and V89) on a Rust crate
written by the same author after it, from the text but sharing comments,
error texts and the generated key format with the Python code. Both implementations have one
author, so their agreement shows that one reading of the specification
is consistent, not that the text is unambiguous; the vectors fix
several choices the text leaves open (Rust README). The criterion is
met only when a bridge written by someone else passes the vectors. That
is the next step. Mapping the contract onto a product runtime
(Termenbox, same author, closed source; as reported by the author on 2026-10-28, a release of the crate is vendored there and the server is not yet wired to it) has already changed the contract: the
resume token held by the client, the single attached client, answer
binding, the coverage limits and the timer derivation were revised as a
result. That is a finding about the draft, not independent
confirmation.

% ================================================================================
